\chapter{Introduction}

\section{Objective}
Develop an open-source tool that computes portfolio weights for Brazilian stocks with the
Black-Litterman (BL) model, incorporating the manager's views. The project has two products:
\begin{itemize}
  \item a web application, in which the user selects assets from a catalog and obtains the weights
        of a portfolio for the current date;
  \item a documented backtest study, which evaluates the method against benchmarks and whose
        results are displayed in the application.
\end{itemize}

\section{Initial scope}
\begin{itemize}
  \item Stocks and units traded on B3.
  \item Asset selection by liquidity, based on the B3 historical quotes files (COTAHIST).
  \item A study portfolio with the 10 most liquid assets at each rebalancing date.
  \item Daily data in the application and weekly data in the study.
  \item One portfolio per run in the application, with at least 5 assets from the catalog.
  \item Risk-free asset: CDI.
  \item Backend deployed on PythonAnywhere and a separate frontend, initially in Streamlit.
\end{itemize}

\section{Target audience}
Investors or portfolio managers with basic knowledge of financial markets who need an economic
method to set the weights of a portfolio of liquid Brazilian stocks.

\section{Definitions}
\begin{description}
  \item[Catalog:] list of the 30 most liquid tickers available for selection in the application.
  \item[Market weights ($w_{mkt}$):] market capitalization weights, renormalized within the universe.
  \item[Prior ($\Pi$):] equilibrium returns implied by the market, consistent with the CAPM.
  \item[View ($P$, $Q$):] the user's opinion on the return of one asset (absolute) or of one asset
        relative to another (relative).
  \item[$\Omega$:] uncertainty matrix of the views.
  \item[$\tau$, $\delta$:] prior uncertainty scalar and risk aversion coefficient.
  \item[Posterior ($\mu_{BL}$):] expected returns that combine the prior and the views.
\end{description}
